\documentclass[11pt]{article}
\usepackage{manual_docs_style}

\begin{document}

\docTitle{Stage: Agentic Rollout}
\phantomsection\label{docstart:agentic_rollout}

This stage runs the agentic experiments. 
These experiments generate artifacts that are later submitted for evaluation or used for additional post-processing. 
It also provides a convenient way to manage, resume, and inspect rollouts.

There are two main scripts and a third ancillary one:
\begin{itemize}
  \item \code{workflows/rollout/question_run_orchestrator.py}, which orchestrates the runs. It supports discriminated unions, a resume option, and the plan/materialize/reconcile pattern.
  \item \code{workflows/rollout/run_single_question.py}, which executes the individual runs in containers.
  \item \code{workflows/rollout/cleanup_tb_question_resources.sh}, which kills any running container, i.e. if we run --keep-container.
\end{itemize}

This stage depends only on the artifact \code{tsENV_questions} export bundle. It relies on \code{terminal_bench} to actually run the experiments.

\section*{Orchestrator}
\begin{DocCode}
python workflows/rollout/question_run_orchestrator.py \
  --tasks-dir tsENV_questions \
  --agent-id AGENT_ID [AGENT_ID ...] \
  [--question-slug QUESTION_SLUG [QUESTION_SLUG ...]] \
  [--row-slug ROW_SLUG [ROW_SLUG ...]] \
  [--shot-slug SHOT_SLUG [SHOT_SLUG ...]] \
  --model SIMULATOR [SIMULATOR ...] \
  [--dry-run] \
  [--resume configuration-file-path [--heal]] \
  [--number-of-runs-in-parallel-per-agent <n>] \
  [--tag <tag>] \ 
  [--call-adapter] \
  [--name <name>]
\end{DocCode}
Note that:
\begin{itemize}
  \item When \code{--resume} is passed, no other argument can be passed. (The number of runs in parallel per agent is read from there).
  \item Only one of \code{--question-slug}, \code{--row-slug},  \code{--shot-slug}, \code{--resume} can be passed simultaneously.
  \item \code{--tag} Tag to attach to each run. Default: DEBUG
  \item \code{--name} Name of the configuration file. If name is not specified, then it's the current <timestamp>. If a file already exists with the same name it throws an error.
\end{itemize}


Internal steps:
\begin{enumerate}
  \item Validate the command-line arguments
  \phantomsection\label{sec:agentic-rollout-run-configurations}
  \item Build a JSON configuration file at \code{run_configurations/<name>.json}:
  \begin{DocCode}
    {
      "timestamp": "...",
      "name": "...",
      "models": ["<simulator_id>", "..."],
      "number_of_runs_in_parallel_per_agent": 1,
      "total_runs": 1,
      "runs": [
        {
          "question_schema": "...",
          "question_slug": "...",
          "model": "...",
          "agent_id": "...",
          "agentic_run_id": "...",
          "status": "PENDING",
          "path_to_the_run": "terminal-bench/runs/<agentic_run_id>",
          "agent_prompt": "...",
          "question": {
            "recipe_info": { "..." : "..." }
          }
        }
      ]
    }
    \end{DocCode}
  where:
  \begin{itemize}
    \item \code{question_schema} is from \code{tsENV_questions/<simulator_id>/questions.json::question}
    \item \code{status} can be one of PENDING|ERROR|DONE|RUNNING
    \item \code{path_to_the_run}: location of the run of the form \code{terminal-bench/runs/<agentic_run_id>}. 
    Note this does not exist if the run is not effectively run.
    \item \code{agent_prompt}: the user prompt that starts the experiment. 
  \end{itemize} 
  This script internally works as follow
    \begin{enumerate}
      \item It resolves any placeholder values with their corresponding values, specifically \code{questions.<question_slug>.question_text.possible_interventions}.
    \end{enumerate}
    \item if \code{--call-adapter} is passed then it also calls \code{terminal-bench/adapters/tsENV/run_adapter.py} (i.e step 2) of the main driver. 
    Note that if the targer folder in \code{terminal-bench/tasks_runtime/manual/<simulator_id>/<question_slug>/question_0/agent_payload} already exists, the script will first delete it and recreate it.
    \item \code{question.recipe_info.*}: All the information stored at \code{question.recipe_info.*}
    \item The agent prompt is created from \hyperref[sec:generate-question-question-text]{\code{question_text}} by calling \code{shared/prompts.py::render_tsenv_agent_prompt}.

  Note instead with \code{--resume} \code{configuration-file-path} we load the file and skip this step. 
  If \code{--heal} is passed we also delete \code{path_to_the_run} and set again to "PENDING" the runs that are set to "ERROR"
  \item Call \code{workflows/rollout/run_single_question.py},   \code{--number-of-runs-in-parallel-per-agent} times in parallel.
  \item Update the status of the configuration files.
  \item When everything has finished run \code{docker image prune} and \code{docker buildx prune}
\end{enumerate}

\section*{Main Driver}
This script actually runs the jobs and creates the output artifacts. It can be called multiple times, provided that the \code{run_id} is different.
Example invocation:
\begin{DocCode}
python workflows/rollout/run_single_question.py \
  --model SIMULATOR \
  --question-slug QUESTION_SLUG \
  --agent-id AGENT_ID \
  --agentic-run-id AGENTIC_RUN_ID \
  [--keep-container] \
  [--rollout-mode SINGLE|ORCHESTRATOR]
\end{DocCode}

Additional flags:
\begin{itemize}
  \item \code{--keep-container}: keeps the container alive after the run.
  \item \code{--rollout-mode}: defines the rollout mode if it's a single run or an orchestrator driven run. 
\end{itemize}

Internal steps:
\begin{enumerate}
  \item Check that all required environment variables are set and validate profiles.
  \item Call \code{terminal-bench/adapters/tsENV/run_adapter.py}, which uses \code{adapter.py} to create the folder uploaded to the container by performing the following steps.
  Note that this step is skipped if the \code{terminal-bench/tasks_runtime/manual/<simulator_id>/<question_slug>} already exists and the \code{question.json::version} matches the one in \code{scenario_info.json::question_schema}
  \begin{enumerate}
    \item Obtain the anonymised time series with its corresponding noise analysis by calling \code{materialize(...)}.
    \item Create the \code{agent_payload} folder with the train and test samples, plus \code{train_labels.json} when needed, and copy the files there in Parquet format.
    \code{train_labels.json} is keyed by the sample UUID filename, i.e. \code{"<uuid>.parquet": "mass"} in case \termDescLevel != none, else it is \code{"<uuid>.parquet": "label_0"} where the possible labels go up to the number of possible classes.
    \item Create the agent user prompt by calling \code{shared/prompts.py::render_tsenv_agent_prompt(question_text: Mapping[str, Any]) -> str}. 
  \end{enumerate}
  The output of this step is \code{terminal-bench/tasks_runtime/manual/<simulator_id>/<question_slug>} and contains among other things:
  \begin{itemize}
    \item \code{task.yaml} with instruction which is exactly the same prompt given to the agent
    \item \code{noise_analysis.json}: a JSON file containing, for each UUID sample in \code{test_samples} and optionally \code{train_samples}, the output of \code{quantify_effect_snr} for the SNR of each noise profile used to materialize that sample (without any reference) 
    \begin{DocCode}
      {
        "17a033a89614421b8bb520c3bb286b60.parquet": {[...]},
        "1b4abf9677a956e4bb90051b4494f4b3.parquet": {[...]}
      }
    \end{DocCode}
  \end{itemize}
  \item Call \verb|uv run tb run ...| to start the agentic run using terminal-bench harness.
  \begin{itemize} 
    \item Copy \code{terminal-bench/tasks_runtime/manual/<simulator_id>/<question_slug>/question_0/agent_payload} to a \code{tmp/<runid>/agent_payload} folder to avoid any race condition
    \item Inside \code{/app} in the container, copy \code{tmp/<runid>/agent_payload}/
    This should be the folders \code{test_samples/<uuid>.parquet}, and optionally if the question has train samples, should also include \code{train_samples/<uuid>.parquet}, and \code{train_labels.json}.
    \item Run the agent, by using the command detailed in Platform Startup Commands in \code{docs/overleaf_paper/src/appendix/appendix_agents.tex}.
    \item Collect all files in the \code{/app} folder, except the original inputs, and move them to \code{artifacts/}.
    If the \code{platform=="opencode"}, also run the OpenCode export procedure:
    \begin{enumerate}
      \item After the agent finishes, export the completed OpenCode session inside the container.
      \item Redirect stdout directly to the mounted agent-log file:
      \begin{DocCode}
opencode export <session_id> > /agent-logs/opencode-session-<session_id>.json
      \end{DocCode}
      \item Do not capture \code{opencode export} stdout through a pipe before writing the file, because OpenCode has had truncated JSON failures when large export payloads are written through piped stdout.
      \item Validate that \code{/agent-logs/opencode-session-<session_id>.json} is valid JSON before accepting the run artifact.
    \end{enumerate}

  \end{itemize}
  \item Cleanup depends on the rollout mode.
       For parallel orchestrated rollout, Terminal-Bench must run \code{docker compose down --volumes} and must not remove images, because all parallel trials for \code{question_0} share the client image tag \code{tb__question_0__client}. 
       Removing images with \code{--rmi all} can race with another trial that has just built the image and is about to start its container. 
       For isolated single-run debugging, \code{docker compose down --rmi all --volumes} is allowed.
  \item Run the post-rollout steps:
  \begin{itemize}
    \item It calls the Artifact scoring script \code{python workflows/rollout/evaluate_artifact.py <agentic_run_id>}
    \item It calls the trajectory export \code{python workflows/trajectory/export_atif_trajectory.py <agentic_run_id>}
    \item It calls the trajectory evaluation \code{python workflows/trajectory/trajectory_evaluation_programmatic.py <agentic_run_id>}
  \end{itemize}
\end{enumerate}
  
Note that:
\begin{itemize}
\item We choose Parquet because it is language-agnostic and platform-agnostic.
\item The container has no internet access except for model API calls routed through a proxy, and Python 3.13 is installed with the ML libraries defined at \code{terminal-bench/adapters/tsENV/template/Dockerfile}.
\end{itemize}

\phantomsection
\section*{Materialized Output}
\label{sec:agentic-rollout-materialized-output}
Stored under:
\begin{DocCode}
terminal-bench/runs/<ID>/question_0
\end{DocCode}

The directory contains:
\begin{itemize}
  \item \code{scenario_info.json}
  \begin{itemize}
    \item identifier related \code{question_id}, \code{agent_id}, \code{tag}
    \item \code{question_schema} from \code{tsENV\_questions/<simulator_id>/questions.json::question}
    \item \code{configuration_file_name}: if it was started using a configuration file name, it reports the name here
    \item instruction related \code{instruction_agent_format}
    \item recipe related \termDescLevel, \code{shot}, \code{noise_level}
    \item \code{agent_started}: It's True if the agent started (i.e. the agent output at least one "agent") message, otherwise it's false
  \end{itemize}
  \item \code{results.json}
  \item \code{artifacts/*}
  \item \code{agent_logs/rollout}
  \item \code{agentic-final-response.json}
\end{itemize}

\subsection*{Kill running resources}

The cleanup command kills any running container, i.e. if we run \code{--keep-container}:
\begin{DocCode}
workflows/rollout/cleanup_tb_question_resources.sh [--run-id RUN_ID ...] [--dry-run]
\end{DocCode}
By default it removes all the running containers and any resource that starts with \code{/question_0-}.
If \code{--dry-run} is passed then it only lists what will be deleted.

\end{document}
